\chapter*{Preface: Two Verbs and a Halt}
\addcontentsline{toc}{chapter}{Preface: Two Verbs and a Halt}

This is the third reading of one machine. Four volumes walk the same fifteen chapters; each writes them in the
language where the machine's face comes out sharp, and English is the fallback, never the medium. The first volume
read the machine as a construction and spoke the logic of building. This one reads it as a \emph{compiler} --- a
program that writes its own source and then reads it back --- and speaks the language such a program is native to:
formal systems, the theory of communication, and the numerical analysis of a finite instrument that can only ever
compute to a floor. There are no forces in this book, no fields, no orbits; where a physical name would be the
natural word, this volume keeps the computational one and hands the physics to the volume that can cash it. What the
other gauges call by their own names, this gauge calls by grammars, codes, and decisions.

Two verbs run under everything here, and they are worth setting down at the door. To \emph{tange} is to select by a
decidable characteristic --- to reach into a collection and pick out \emph{this one, not that one}, because a finite
procedure, run to a halt, reports them different. To \emph{funge} is the opposite gesture: to pool together, under
one tag, everything a finite procedure cannot tell apart. Tange divides by a decided difference; funge unites by a
decided sameness. Neither verb reaches past what the machine can already settle in finite time. A tange is a
decidable selection; a funge is a decidable grouping; and between the two of them, iterated, the machine builds
everything it will ever build. The reader who keeps both verbs in hand --- and remembers that each is bounded by a
procedure that must \emph{halt} --- will find the whole construction easier to walk.

They are also, this volume argues, the whole of what it means for a compiler to \emph{measure}. To measure is first to
tange --- to decide a symbol apart from every symbol it is not --- and then to funge --- to declare it one with
everything no decision can separate from it. This is not a metaphor borrowed from elsewhere; it is the plainest fact
about a symbol, older than any computer. A symbol means nothing in isolation and means only by its contrast with the
symbols it is not; Shannon~\cite{shannon1948} made that quantitative, a mark carrying information exactly to the
degree it could have been another. The machine takes the old idea literally: its symbols \emph{are} their
distinguishability, a mark defined by the finite test that parts it from the rest. The test is the tange; the pooling
of the indistinguishable is the funge; and every reading in these pages is built from those two acts and a procedure
that halts.

There is a place where the first verb stops being an ordinary word and becomes the deepest fact about computation, and
this book walks up to it on purpose and then declines to cross. Tange one symbol from a finite alphabet and no one
objects; a finite procedure decides it and halts. But ask to tange, in general, whether an arbitrary procedure will
ever halt --- to decide, for every program, the one bit that says \emph{stops} or \emph{runs forever} --- and you have
asked for something no machine can deliver. Turing~\cite{turing1936} proved it: there is no decision procedure for
halting, no finite test that settles the question for all cases. G\"odel~\cite{godel1931} had shown the same wall from
the other side --- that any system strong enough to describe its own arithmetic contains true statements it cannot
prove --- and Church~\cite{church1936} a third time, that the plainest questions about definability admit no algorithm.
The undecidable is not a gap in anyone's cleverness; it is a theorem, a standing warning about how much a finite
instrument can and cannot settle. To reach past it --- to help oneself to the halting bit as though it were free ---
is to invoke an \emph{oracle}, a selection the machine cannot exhibit, tange promoted past the finite and granted for
nothing.

This book takes a side, and the side is the finite one. The machine at the center of these pages tanges and funges
without pause --- it does nothing else --- but it only ever tanges by procedures that halt, and only ever from
alphabets it can genuinely decide. It calls no oracle; it never helps itself to a bit it cannot compute; it never
posits a distinction it cannot exhibit in finite time. When it names a thing, it does so by counting the thing into
one of finitely many bins, and a decision among finitely many bins needs no oracle at all --- you can just run it and
watch it halt. The consequence is exact and checkable: audit what the machine's results actually rest on and no oracle
is among them; the only thing it ever assumes is that a halting procedure halts. The book does not avoid the oracle as
a matter of taste. It earns every selection one finite, halting decision at a time, and it can show that it did.

That discipline is the reason the book can end where it does. There is a fence in the theory of information that is the
exact computational form of this refusal, and it belongs here at the door because the whole ending leans on it. A
finite string carries only so much information; most strings are \emph{incompressible}, holding no shorter description
than themselves, and Chaitin~\cite{chaitin1974} made the boundary sharp: there are numbers a finite machine provably
cannot generate, magnitudes whose every digit is beyond the reach of any short program. A machine that helped itself
to an oracle could posit any such number it pleased and dress the wish as a result --- it could hand you the exact
value of the constant this book reaches for and dare you to object. This machine cannot, and will not. It tanges only
what it can decide and funges only what no decision separates; and when it runs clean out of decidable distinctions ---
when its bisection reaches the floor below which its own resolution cannot tell two candidates apart --- it
\emph{halts}, and hands back not a fabricated point but an honest interval. The width of that interval is the exact
measure of the information its finite, halting procedures could not resolve. The two verbs are the method; the oracle
is what the first becomes when it is handed the infinite for free; and the whole of this book's honesty is its refusal
of that gift. It tanges, it funges, and it halts --- and because it halts, the answer it finally gives is a thing it
earned rather than a thing it wished for.
